\section{Consequence I: the cache is editable}
\label{sec:editable}

\Cref{fig:ops} contrasts the two cache operations this section and the next make precise: editing appends
a salient correction; composing repositions and splices precompiled KV.

\paragraph{Naive editing fails---by design.} The mechanism predicts the in-place edit will fail, and it
does: refreshing the field's KV while reusing the stale downstream leaves the decision at the old value
(\Cref{fig:teaser}b). The prefix before the field is genuinely reusable (deviation $0.0$); the problem is
that the conclusion was already memoized after it.

\begin{figure}[t]
\centering
\includegraphics[width=\linewidth]{fig_operations}
\caption{\textbf{Two cache operations.} (a) \emph{Editable}: append a salient erratum that amends the
stale notes ($O(1)$ extra work, the whole prefix is reused). (b) \emph{Composable}: precompute a skill's
KV in isolation, RoPE-reposition it to the target positions, and splice it in ($O(L)$, no reprefill).}
\label{fig:ops}
\end{figure}

\paragraph{The robust fix is an erratum, not recomputation.} If the stale notes carry an old conclusion,
the cheapest correct intervention is to append a \emph{salient erratum} that supplies the corrected state
late in the context---``\texttt{[STATE UPDATE] field $\rightarrow$ new; overrides any earlier value and
conclusion}''---so the decision token attends to a fresh, authoritative note (verbatim template in
\Cref{app:examples}). Appended after the original
field (\texttt{field+erratum}), this matches the strong \emph{hoist-to-end} baseline's oracle correctness
\emph{without rewriting the prompt}: under chain-of-thought the erratum recovers the flipped decision at
$0.97$ on Qwen3-8B (\Cref{fig:edit}a), and the edit is append-only, so it composes with prefix caching
(\Cref{sec:systems}).

\paragraph{When the ${\sim}1\%$-compute edit alone suffices: reasoning-gated and model-dependent.}
Surprisingly, under explicit reasoning the cheap in-place edit \emph{by itself} can recover the decision,
because the chain re-reads the (now-refreshed) field and re-derives the conclusion rather than trusting
the stale notes. This recovery is strongly model-dependent (\Cref{fig:edit}b): how ``sticky'' the memoized
conclusion is varies non-monotonically across model sizes rather than tracking scale, so the minimal
amount of selective recomputation needed, \fieldsel{}@$K$ (refresh the field plus the $K$ highest-effect
downstream tokens), varies from $K^\star\approx 4$ at $8$B to $>64$ at $4$B (\Cref{fig:edit}c). Without reasoning, no small
$K$ helps ($0.00$ recovery at every scale). \fieldsel{}@$K$ is therefore a genuine but
\emph{unreliable} surgical tool---effective when the conclusion is not sticky, model-dependent otherwise;
we report it honestly rather than as a default.

\paragraph{A frontier, not a winner.} A controlled comparison (full table in \Cref{app:edit}) shows no single
dominant method. Hoist-to-end is cheapest but demands prompt surgery and pre-identification of every field;
\texttt{field+erratum} matches it with no surgery at a one-line append; \inplace{} is near-free but only
under reasoning; a KV-deviation-ranked selective recompute (CacheBlend-style) underperforms here because it
chases changed keys rather than the tokens that \emph{memoized the conclusion}. The practical
recommendation is \texttt{field+erratum} as the robust default, with \inplace{} as a free fast-path for
reasoning models.

\paragraph{Why not edit the weights?} A natural objection: to act on a changed field, why not edit
the model's \emph{weights} (ROME/MEMIT) or fine-tune, rather than the cache? We compare on the paper's
gated task (status \texttt{pending}$\rightarrow$\texttt{shipped}, so \emph{cancel}$\rightarrow$\emph{deny})
against a faithful rank-one ROME---validated on the canonical factual edit first (``the Eiffel Tower is
in'' \emph{Paris}$\rightarrow$\emph{Rome}, locality intact), so the baseline is not crippled---and a LoRA
fine-tune (\Cref{tab:weight}; methodology in \Cref{app:systems}). All of \fielderr{}, ROME, and LoRA \emph{succeed} at flipping the target
decision. But a weight edit is \emph{global}: the same model instance can no longer hold
\texttt{status=shipped} for one request and \texttt{pending} for another, so \emph{all} concurrent orders
that are genuinely still pending are wrongly flipped (cross-request contamination $1.0$), and half of an
unrelated decision battery drifts ($0.5$)---the ROME/fine-tune specificity tax---at $3$--$6$\,s per edit
(plus a one-time covariance pass for ROME). The append-only erratum lives in a \emph{per-sequence} cache:
zero cross-request contamination, zero collateral, $114$\,ms, and it composes with prefix caching
(\Cref{sec:systems}). Weight editing targets \emph{durable, global facts}; mutable per-request, per-turn
state is the wrong job for it---which is precisely the niche the editable KV cache fills.

\begin{table}[t]
\centering\small
\caption{\textbf{KV editing vs.\ weight editing} for mutable per-request state (Llama-3.1-8B). Weight edits
succeed at the target yet are global (contaminate concurrent requests), damage unrelated decisions, and are
$30$--$50\times$ slower per edit.}
\label{tab:weight}
\begin{tabular}{lcccc}
\toprule
method & flips decision? & edit latency & cross-req.\ contamination & collateral \\
\midrule
\fielderr{} (KV) & \ding{51} & $114$\,ms & $\mathbf{0}$ & $\mathbf{0}$ \\
\inplace{} (KV) & \ding{55}\,(no CoT) & $71$\,ms & $0$ & $0$ \\
ROME (rank-one) & \ding{51} & $5.6$\,s\,$+\,11.6$\,s cov & $1.0$ & $0.5$ \\
LoRA fine-tune & \ding{51} & $3.1$\,s & $1.0$ & $0.5$ \\
\bottomrule
\end{tabular}
\end{table}

\begin{figure}[t]
\centering
\includegraphics[width=\linewidth]{fig3_editable}
\caption{\textbf{Editing the cache.} (a) Editing the field's KV is ignored; \texttt{field+erratum} reaches
the hoist-to-end oracle under CoT. (b) The stickiness of the memoized conclusion is model-dependent, so
field-only reasoning recovery varies non-monotonically across model sizes. (c) \texttt{field+selective}@$K$: the
minimal recompute $K^\star$ to reach full quality is model-dependent.}
\label{fig:edit}
\end{figure}
